\documentclass[10pt]{article}
\usepackage[a4paper,margin=17mm]{geometry}
\usepackage{amsmath,amssymb,booktabs,array,microtype}
\usepackage[colorlinks=true,linkcolor=black,urlcolor=blue]{hyperref}
\usepackage{enumitem}
\setlist{nosep,leftmargin=5mm}
\setlength{\parindent}{0pt}
\setlength{\parskip}{4pt}
\newcommand{\Cbar}{\overline{\mathbf C}}
\newcommand{\Dt}{\widetilde{\mathbf D}}

\begin{document}
\begin{center}
  {\Large\bfseries Inter-arc continuity mathematics: before and after review}\\[2pt]
  {\small tudatpy PR 828 review follow-up, 22 July 2026}
\end{center}

\textbf{Quantities used by both implementations.}
For body $b$ and constrained arc pair $k=(\ell,r)$, evaluated at connection epoch $t_{c,bk}$,
\begin{align*}
 \mathbf d_{bk} &:= \mathbf x_{r,b}(t_{c,bk})-\mathbf x_{\ell,b}(t_{c,bk}),\\
 \mathbf D_{bk} &:= \frac{\partial\mathbf d_{bk}}{\partial\mathbf p}
 =\mathbf M_{r,b}(t_{c,bk})-\mathbf M_{\ell,b}(t_{c,bk}),\\
 \Dt_{bk} &:= \mathbf D_{bk}\,\operatorname{diag}(\mathbf s)^{-1}.
\end{align*}
Here $\mathbf p$ is the estimated parameter vector, $\mathbf M$ denotes the relevant rows of the full
state-transition/sensitivity matrix, and $\mathbf s$ contains Tudat's observation-design-matrix column
normalization factors.  Let $m$ be the number of scalar observation residuals and
\[
 \Cbar_{bk}:=\tfrac12(\mathbf C_{bk}+\mathbf C_{bk}^{T}),\qquad
 m_d:=\sum_{b,k}\operatorname{rank}(\Cbar_{bk}).
\]
$\mathbf C_{bk}$ is the \emph{continuity} weight supplied through \nolinkurl{weight_matrices}, or the diagonal
matrix constructed from \nolinkurl{position_weights} and \nolinkurl{velocity_weights}.  It is distinct from the
observation weight matrix $\mathbf W_o$.  The symmetric part matters only for input asymmetry within the
accepted numerical tolerance.

\textbf{Before the mathematical correction.}
The code used
\[
 \mathbf W^-_{bk}=\frac{1}{\mu_s m_d}\mathbf C_{bk},
\]
where $\mu_s$ is the settings object's \texttt{constraint\_scaling\_factor}.  Ignoring round-off-level
asymmetry for this paragraph, the normal-equation additions were
\begin{align*}
 \Delta\widetilde{\mathbf A}^{-}
 &=\sum_{b,k}\Dt_{bk}^{T}\mathbf W^-_{bk}\Dt_{bk},&
 \Delta\widetilde{\mathbf q}^{-}
 &=-\sum_{b,k}\Dt_{bk}^{T}\mathbf W^-_{bk}\mathbf d_{bk}.
\end{align*}
Consequently, the solve represented the local quadratic objective
\[
 J^-_{\rm solve}=\tfrac12\mathbf r^T\mathbf W_o\mathbf r
 +\tfrac12\sum_{b,k}\mathbf d_{bk}^{T}\mathbf W^-_{bk}\mathbf d_{bk}.
\]
However, iteration selection added the reported constraint cost without the factor $1/2$:
\[
 J^-_{\rm select}=\tfrac12\mathbf r^T\mathbf W_o\mathbf r
 +\sum_{b,k}\mathbf d_{bk}^{T}\mathbf W^-_{bk}\mathbf d_{bk}.
\]
Thus selection evaluated a different observation/constraint balance from the one used by the correction step.
The constraint also lacked the factor $m$ required to retain the relative normalization of the averaged
observation and discontinuity terms in the cited Orbit14 model.

For an accepted slightly non-symmetric $\mathbf C$, the old implementation had an additional inconsistency:
the explicitly symmetrized normal matrix and scalar quadratic cost used $\Cbar$, while the right-hand side used
raw $\mathbf C$.

\newpage
\textbf{After the mathematical correction.}
Every continuity term now uses
\[
 \boxed{\mathbf W^+_{bk}=\frac{m}{\mu_s m_d}\Cbar_{bk}}.
\]
The single implemented objective, its linear term, and its curvature are
\begin{align*}
 J^+ &= \tfrac12\mathbf r^T\mathbf W_o\mathbf r
       +\tfrac12\sum_{b,k}\mathbf d_{bk}^{T}\mathbf W^+_{bk}\mathbf d_{bk},\\
 \Delta\widetilde{\mathbf A}^{+}
 &=\sum_{b,k}\Dt_{bk}^{T}\mathbf W^+_{bk}\Dt_{bk},&
 \Delta\widetilde{\mathbf q}^{+}
 &=-\sum_{b,k}\Dt_{bk}^{T}\mathbf W^+_{bk}\mathbf d_{bk}.
\end{align*}
With $\mathbf d(\mathbf p+\delta\mathbf p)\simeq\mathbf d+\mathbf D\delta\mathbf p$, these are exactly the
Hessian/Gauss--Newton term and negative gradient contribution of the same quadratic objective.  The reported
constraint cost is now the second term of $J^+$, so best-iteration selection and the solve use the same model.

The Orbit14 equations used as the reference are
\[
 Q_{\rm ref}=\frac{1}{m}\mathbf r^T\mathbf W_o\mathbf r
 +\sum_{b,k}\frac{1}{\mu_s m_d}\mathbf d_{bk}^{T}\Cbar_{bk}\mathbf d_{bk}.
\]
Therefore
\[
 \boxed{J^+=\frac{m}{2}Q_{\rm ref}},
\]
which has the same minimizer as $Q_{\rm ref}$ and matches Tudat's existing
$\tfrac12\mathbf r^T\mathbf W_o\mathbf r$ convention.  A separate $\mu_s$ per settings object is the
implementation's direct extension of the paper's single penalty parameter.

\textbf{What did not change.}
\begin{itemize}
 \item The discrepancy remains right-arc state minus left-arc state.
 \item Its partial remains right-arc variational rows minus left-arc variational rows.
 \item The estimation right-hand-side sign remains negative for the continuity term.
 \item Component masks and dense positive-semidefinite weights remain supported; $m_d$ is their total rank,
       generalizing $6(n-1)$.
\end{itemize}

\textbf{References.}
G.~Lari et al., ``Orbit determination methods for interplanetary missions: the Orbit14 software,''
\emph{Experimental Astronomy} (2021), Eqs.~(4) and (28),
\href{https://doi.org/10.1007/s10686-021-09823-8}{doi:10.1007/s10686-021-09823-8}.
For comparison, E.~M.~Alessi et al., ``Orbit determination strategy and accuracy for the BepiColombo radio
science experiment,'' \emph{Advances in Space Research} (2012), Eq.~(5), uses a different common-denominator
normalization; the implementation follows the later Orbit14 formulation above.

\end{document}
